% ══════════════════════════════════════════════════════════════
%  EXAMEN 1 — SEGURIDAD QUÍMICA
% ══════════════════════════════════════════════════════════════

\phantomsection
\addcontentsline{toc}{section}{1. Riesgo vs. Peligro}
\section*{\centering\fontsize{12}{18}\selectfont{\textcolor{blueunam}{\textbf{1. RIESGO VS. PELIGRO}}}}

\textit{Explica brevemente la diferencia fundamental entre ``riesgo'' y ``peligro''.}

\vspace{0.3cm}
Según el SGA (Libro Violeta, ONU, 2021), el \textbf{peligro} es la propiedad intrínseca de un agente o situación con potencial de causar daño a las personas, la propiedad o el medio ambiente. Es una característica inherente e invariable de la sustancia. El \textbf{riesgo} es la probabilidad de que ese daño se materialice bajo condiciones específicas de uso o exposición; puede reducirse mediante controles sin que el peligro desaparezca.

% ──────────────────────────────────────────────────────────────

\phantomsection
\addcontentsline{toc}{section}{2. Jerarquía de Controles}
\section*{\centering\fontsize{12}{18}\selectfont{\textcolor{blueunam}{\textbf{2. JERARQUÍA DE CONTROLES}}}}

\textit{Señala los tres primeros niveles de la Jerarquía de Controles.}

\vspace{0.3cm}
\begin{enumerate}
    \item \textbf{Eliminación:} suprimir completamente el peligro en la fuente.
    \item \textbf{Sustitución:} reemplazar el agente peligroso por uno de menor riesgo.
    \item \textbf{Controles de ingeniería:} aislar o contener el peligro mediante medidas físicas (campanas de extracción, ventilación localizada, encerramientos).
\end{enumerate}

% ──────────────────────────────────────────────────────────────

\phantomsection
\addcontentsline{toc}{section}{3. Siglas GHS}
\section*{\centering\fontsize{12}{18}\selectfont{\textcolor{blueunam}{\textbf{3. SIGLAS GHS}}}}

\textit{Explica el significado de las siglas GHS.}

\vspace{0.3cm}
\textbf{GHS:} \textit{Globally Harmonized System of Classification and Labelling of Chemicals} --- Sistema Globalmente Armonizado

% ──────────────────────────────────────────────────────────────

\phantomsection
\addcontentsline{toc}{section}{4. Pictograma de toxicidad}
\section*{\centering\fontsize{12}{18}\selectfont{\textcolor{blueunam}{\textbf{4. PICTOGRAMA DE TOXICIDAD}}}}

\textit{Describe el pictograma relacionado con la toxicidad de un producto químico.}

\vspace{0.3cm}
El SGA define tres pictogramas asociados a toxicidad:

\begin{itemize}
    \item \textbf{GHS06 (calavera con tibias cruzadas):} toxicidad aguda categorías 1, 2 y 3 por inhalación, ingestión o contacto cutáneo. Indica sustancias mortales o muy tóxicas en dosis bajas.

    \begin{center}
        \includegraphics[width=2cm]{../GHS/GHS06.png}
    \end{center}

    \item \textbf{GHS07 (signo de exclamación):} toxicidad aguda categoría 4, irritación cutánea u ocular, sensibilización cutánea, efectos narcóticos, toxicidad para órganos diana por exposición única. También cubre peligro para la capa de ozono.

    \begin{center}
        \includegraphics[width=2cm]{../GHS/GHS07.png}
    \end{center}

    \item \textbf{GHS08 (riesgo para la salud):} peligros crónicos: carcinogenicidad, mutagenicidad en células germinales, toxicidad reproductiva, sensibilización respiratoria, toxicidad para órganos diana por exposición repetida, peligro por aspiración.

    \begin{center}
        \includegraphics[width=2cm]{../GHS/GHS08.png}
    \end{center}
\end{itemize}

% ──────────────────────────────────────────────────────────────

\phantomsection
\addcontentsline{toc}{section}{5. Responsabilidades de los usuarios}
\section*{\centering\fontsize{12}{18}\selectfont{\textcolor{blueunam}{\textbf{5. RESPONSABILIDADES DE LOS USUARIOS}}}}

\textit{Señala las responsabilidades de los usuarios en su interacción con productos químicos.}

\vspace{0.3cm}
Conforme al SGA y la NOM-018-STPS-2015:

\begin{itemize}
    \item Leer y comprender la etiqueta y la hoja de datos de seguridad antes de manipular cualquier sustancia.
    \item Utilizar el equipo de protección personal indicado en la sección 8 de la SDS.
    \item Almacenar las sustancias en las condiciones especificadas y separar materiales incompatibles.
    \item No transvasar a recipientes sin etiqueta ni reutilizar envases originales para otras sustancias.
    \item Reportar de inmediato cualquier accidente, derrame o condición insegura al responsable del área.
    \item Participar en las capacitaciones de seguridad química obligatorias.
    \item No modificar, alterar ni retirar etiquetas de los recipientes.
\end{itemize}

% ──────────────────────────────────────────────────────────────

\phantomsection
\addcontentsline{toc}{section}{6. Contenido mínimo de una hoja de seguridad}
\section*{\centering\fontsize{12}{18}\selectfont{\textcolor{blueunam}{\textbf{6. CONTENIDO MÍNIMO DE UNA HOJA DE SEGURIDAD}}}}

\textit{Describe el contenido mínimo de una hoja de seguridad.}

\vspace{0.3cm}
El SGA (Libro Violeta, cap. 4) establece que toda SDS debe contener \textbf{16 secciones}. A continuación se describe brevemente cada una:

\begin{enumerate}
    \item \textbf{Identificación del producto y del proveedor:} nombre comercial, uso recomendado, datos de contacto del fabricante o distribuidor.
    \item \textbf{Identificación de peligros:} clasificación SGA, pictogramas, palabra de advertencia, frases H y P.
    \item \textbf{Composición e información sobre los ingredientes:} sustancias, mezclas, impurezas relevantes y número CAS.
    \item \textbf{Primeros auxilios:} medidas según vía de exposición (inhalación, ingestión, piel, ojos).
    \item \textbf{Medidas de lucha contra incendios:} agentes extintores adecuados, peligros específicos de combustión, EPP para bomberos.
    \item \textbf{Medidas en caso de vertido accidental:} precauciones personales, ambientales y métodos de contención y limpieza.
    \item \textbf{Manipulación y almacenamiento:} condiciones seguras de uso, incompatibilidades, temperatura y humedad de almacenamiento.
    \item \textbf{Controles de exposición y protección personal:} valores límite de exposición, EPP requerido, controles de ingeniería.
    \item \textbf{Propiedades físicas y químicas:} estado físico, color, olor, pH, punto de fusión/ebullición, inflamabilidad, solubilidad, entre otros.
    \item \textbf{Estabilidad y reactividad:} condiciones a evitar, materiales incompatibles, productos de descomposición peligrosos.
    \item \textbf{Información toxicológica:} vías de exposición, efectos agudos y crónicos, DL\textsubscript{50}, CL\textsubscript{50}.
    \item \textbf{Información ecológica:} toxicidad acuática, persistencia, bioacumulación, movilidad en suelo.
    \item \textbf{Consideraciones relativas a la eliminación:} métodos seguros de disposición de residuos conforme a la normativa aplicable.
    \item \textbf{Información relativa al transporte:} número ONU, clase de peligro, grupo de embalaje, nombre de expedición.
    \item \textbf{Información sobre la reglamentación:} normativa específica de seguridad, salud y medio ambiente aplicable a la sustancia.
    \item \textbf{Otras informaciones:} fecha de elaboración o última revisión, fuentes de datos.
\end{enumerate}

% ──────────────────────────────────────────────────────────────

\phantomsection
\addcontentsline{toc}{section}{7. Etiqueta para producto corrosivo, mutagénico, inflamable y tóxico}
\section*{\centering\fontsize{12}{18}\selectfont{\textcolor{blueunam}{\textbf{7. ETIQUETA: PRODUCTO CORROSIVO, MUTAGÉNICO, INFLAMABLE Y TÓXICO}}}}

\textit{Diseña una etiqueta para un producto corrosivo, mutagénico, inflamable y tóxico.}

\vspace{0.3cm}
Conforme al SGA, los elementos obligatorios de una etiqueta son: identificador del producto, pictogramas, palabra de advertencia, indicaciones de peligro (frases H), consejos de prudencia (frases P) e información del proveedor.

Para este producto los pictogramas son:

\begin{center}
\begin{tabular}{cccc}
    \includegraphics[width=2cm]{../GHS/GHS02.png} &
    \includegraphics[width=2cm]{../GHS/GHS05.png} &
    \includegraphics[width=2cm]{../GHS/GHS06.png} &
    \includegraphics[width=2cm]{../GHS/GHS08.png} \\[4pt]
    {\small Inflamable} & {\small Corrosivo} & {\small Tóxico agudo} & {\small Mutagénico}
\end{tabular}
\end{center}

\vspace{0.3cm}
Palabra de advertencia: \textbf{PELIGRO} (aplica cuando al menos un peligro pertenece a categoría grave). Las frases H y P se seleccionan según la clasificación específica de cada peligro.

% ──────────────────────────────────────────────────────────────

\phantomsection
\addcontentsline{toc}{section}{8. Pictogramas SGA vs. transporte}
\section*{\centering\fontsize{12}{18}\selectfont{\textcolor{blueunam}{\textbf{8. PICTOGRAMAS SGA VS. PICTOGRAMAS DE TRANSPORTE}}}}

\textit{Indica cuál es la diferencia fundamental entre los pictogramas SGA y los pictogramas empleados en el transporte de materiales peligrosos.}

\vspace{0.3cm}
\textbf{Pictogramas SGA:} rombo con borde rojo sobre fondo blanco. Clasifican los peligros intrínsecos de la sustancia para comunicación en el lugar de trabajo y la cadena de suministro, según las clases de peligro del Libro Violeta.

\textbf{Pictogramas de transporte (Recomendaciones ONU / ADR / IATA):} rombo con fondo de color sólido fijo por clase y número de clase ONU visible. Identifican la clase de riesgo durante el tránsito (Clase 3: líquidos inflamables, Clase 6.1: tóxicos, etc.) y están diseñados para ser reconocibles a distancia en vehículos y embalajes.

\textbf{Diferencia fundamental:} el SGA regula la comunicación de peligros en el uso y comercio de la sustancia; el sistema de transporte regula la identificación del riesgo durante el traslado físico. Sus criterios de clasificación, códigos de color y formatos son independientes y no equivalentes entre sí.

% ──────────────────────────────────────────────────────────────

\phantomsection
\addcontentsline{toc}{section}{9. Fuente de información del SGA}
\section*{\centering\fontsize{12}{18}\selectfont{\textcolor{blueunam}{\textbf{9. FUENTE DE INFORMACIÓN DEL SGA}}}}

\textit{¿Cuál es la fuente de información para el SGA y sus pictogramas?}

\vspace{0.3cm}
La fuente oficial es el documento \textbf{ST/SG/AC.10/30} de las Naciones Unidas, conocido como \textbf{Libro Violeta} (\textit{The Globally Harmonized System of Classification and Labelling of Chemicals}). En México se implementa mediante la \textbf{NOM-018-STPS-2015}.

% ──────────────────────────────────────────────────────────────

\phantomsection
\addcontentsline{toc}{section}{10. Información mínima de una hoja de seguridad}
\section*{\centering\fontsize{12}{18}\selectfont{\textcolor{blueunam}{\textbf{10. INFORMACIÓN MÍNIMA DE UNA HOJA DE SEGURIDAD}}}}

\textit{¿Cuál es la información mínima que debe contener una hoja de seguridad?}

\vspace{0.3cm}
Toda SDS debe contener las \textbf{16 secciones} establecidas por el SGA (véase pregunta 6 para el detalle de cada una). Como mínimo indispensable deben estar presentes: identificación del producto y peligros, composición, primeros auxilios, manejo seguro, controles de exposición y EPP, propiedades fisicoquímicas, estabilidad, toxicología e información del proveedor con fecha de revisión.

% ──────────────────────────────────────────────────────────────

\phantomsection
\addcontentsline{toc}{section}{11. Características de la NOM-018-STPS}
\section*{\centering\fontsize{12}{18}\selectfont{\textcolor{blueunam}{\textbf{11. CARACTERÍSTICAS DE LA NOM-018-STPS}}}}

\textit{Indica tres características de la NOM-018-STPS.}

\vspace{0.3cm}
\begin{enumerate}
    \item Establece el sistema de identificación y comunicación de peligros y riesgos de las sustancias químicas peligrosas en los centros de trabajo de México, armonizado con el SGA de la ONU.
    \item Obliga al uso de hojas de datos de seguridad de 16 secciones en idioma español y etiquetas con los elementos SGA para toda sustancia química peligrosa presente en el centro de trabajo.
    \item Define obligaciones diferenciadas: el patrón debe elaborar o gestionar las SDS, capacitar periódicamente a los trabajadores y mantener un inventario actualizado de sustancias peligrosas; los trabajadores deben conocer, consultar y aplicar la información de etiquetas y SDS.
\end{enumerate}